% 第一章 绪论

\chapter{绪论}

\section{研究背景与意义}

数字音乐平台在音乐传播、作品分发和用户消费方面发挥了重要作用，但传统平台通常以中心化服务为核心，创作者的发行流程、作品访问控制、收益分配和交易记录均依赖平台内部系统完成。对于独立音乐人和小型创作团队而言，这种模式可能带来发行门槛较高、收益分配不透明、版权凭证难以外部验证等问题。

区块链和IPFS为音乐资源管理提供了新的技术路径。区块链适合记录作品发布、访问凭证、购买交易和收益分账规则，IPFS适合保存音频、封面和metadata等内容资源。二者结合后，可以形成链上身份与链下内容寻址协同的音乐资源管理方式。与此同时，桌面端音乐应用能够提供更稳定的播放体验、本地缓存、下载管理和系统级交互能力，适合作为链上音乐资源的消费入口。

因此，本文围绕FreeFlow项目，设计并实现一个基于区块链与IPFS的链上音乐发行与播放系统。该系统以纯系统开发为主线，重点实现创作者发布、链上访问控制、用户购买播放、评论互动和链上资源检索等功能，并通过链上音乐资源检索与排序模块增强系统的资源发现能力。

\section{国内外研究现状}

本节后续应结合正式参考文献进行扩写，建议从以下三个方向展开。

\subsection{数字音乐发行系统研究现状}

现有数字音乐平台通常采用中心化的版权管理、分发和收益结算方式。相关研究主要关注数字版权管理、流媒体播放、音乐推荐和创作者平台建设等内容。本文后续应补充国内外主流音乐平台、数字版权管理和创作者经济方面的研究。

\subsection{区块链与IPFS在数字内容管理中的应用}

区块链具有不可篡改、可追溯和可编程合约等特性，常用于数字资产凭证、版权确权和交易记录。IPFS采用内容寻址方式保存文件，可将资源CID作为链上metadata的一部分。本文后续应补充NFT、ERC-1155、去中心化存储和Web3内容分发方面的研究。

\subsection{音乐检索与资源排序研究现状}

音乐检索通常涉及关键词匹配、元数据检索、内容特征检索和个性化推荐等方向。FreeFlow当前处于链上音乐资源冷启动阶段，用户行为数据尚不充分，因此本文选择可解释的检索排序模型作为系统算法模块。后续应补充信息检索、搜索排序和音乐资源检索相关文献。

\section{研究内容与主要工作}

本文的主要工作包括以下几个方面：

\begin{enumerate}
  \item 设计FreeFlow链上音乐发行与播放系统的总体架构，明确桌面端、Web2.5后端、智能合约和IPFS存储之间的职责边界。
  \item 设计并实现创作者工作台，支持作品草稿管理、音乐信息编辑、素材上传、metadata生成、链上发布和访问状态查询。
  \item 设计并实现Web2.5后端服务，提供SIWE登录、Pinata上传代理、发行状态管理、链上资源索引、购买记录、评论和用户资料等能力。
  \item 设计并实现智能合约模块，通过ERC-1155访问凭证、平台发布合约和收益分账合约完成链上发布、购买授权和收入分配。
  \item 设计并实现链上音乐资源检索与排序模块，综合文本字段、链上身份字段、新鲜度和热度信号完成资源排序。
  \item 设计并实现桌面端播放、音频代理、链上音乐库、评论和下载等用户侧功能，形成从创作者发布到用户播放的业务闭环。
\end{enumerate}

\section{论文组织结构}

本文共分为八章，各章安排如下。

第一章为绪论，介绍研究背景、研究现状、主要工作和论文组织结构。

第二章介绍系统使用的相关技术，包括Electron、React、Node.js、Express、Prisma、SIWE、IPFS、Solidity、ERC-1155以及音频代理和Range请求等。

第三章进行系统需求分析，说明系统应用场景、用户角色、功能需求、非功能需求和可行性。

第四章进行系统总体设计，介绍系统总体架构、功能模块、业务流程、数据库设计、智能合约设计和安全设计。

第五章介绍系统核心模块实现，包括Profile登录、创作者工作台、素材上传、链上发布、购买验证、播放器和评论模块。

第六章介绍链上音乐资源检索与排序模块，说明索引数据、查询归一化、候选召回、特征评分、排序公式和评估方式。

第七章进行系统测试与效果分析，包括功能测试、接口测试、合约测试、检索排序测试和性能测试。

第八章总结全文工作，分析当前不足，并提出后续改进方向。
